% Manuscript skeleton (Elsevier elsarticle, target: Computers & Security).
% Rules: every number in Sections 6-7 comes from \input{tables/...} or figures/... generated by
% scripts/make_tables.py / make_figures.py. Text marked \TODO{} is written only after results exist.
\documentclass[review,12pt]{elsarticle}
\usepackage{amsmath,amssymb,amsthm,booktabs,graphicx,hyperref,xcolor,algorithm,algpseudocode}
\newcommand{\TODO}[1]{\textcolor{red}{[TODO: #1]}}
\newtheorem{proposition}{Proposition}
\DeclareMathOperator{\tr}{tr}
\journal{Computers \& Security}

\begin{document}
\begin{frontmatter}
\title{Generative AI, Agentic AI and RAG: Emerging Security Threats and Defense Mechanisms}
\author{Sagar Korde\TODO{affiliation, co-authors}}
\begin{abstract}
Generative AI systems are no longer single models: they retrieve documents, call tools and act over many
steps. Each addition opens a new attack surface, and defenses built for one surface rarely transfer to the
next. This paper organizes the threat landscape of generative, retrieval-augmented and agentic AI into a
taxonomy of twelve threat classes mapped to pipeline stage and security property. It then proposes the
Wasserstein Trust Score (WTS), which treats every untrusted text unit (a user prompt, a retrieved chunk, a
tool output, an agent plan step) as an empirical distribution over token representations and measures its
Wasserstein-2 distance from a benign reference. A closed-form Bures--Wasserstein estimator makes the score
cheap, and split conformal prediction turns it into a decision rule whose false-positive rate is bounded
without distributional assumptions. \TODO{one sentence of headline results with intervals, from results/verdicts.md}
\end{abstract}
\begin{keyword}
prompt injection \sep retrieval-augmented generation \sep LLM agents \sep optimal transport \sep
Wasserstein distance \sep conformal prediction
\end{keyword}
\end{frontmatter}

% ============================================================================================
\section{Introduction}\label{sec:intro}
\TODO{900 words. Draft after results; contributions list below is fixed by the plan.}
Contributions:
\begin{enumerate}
\item A taxonomy of twelve threat classes across generative, retrieval-augmented and agentic AI, mapped to
pipeline stage and security property (Table~\ref{tab:taxonomy}).
\item The Wasserstein Trust Score: a distribution-level, closed-form optimal-transport detector applied
without modification to prompts, retrieved chunks, tool outputs and agent plan steps, including a
measurement of agent intent drift.
\item Conformal calibration of the score with a finite-sample false-alarm guarantee, validated in
distribution and stress-tested under corpus shift.
\item An evaluation on open and API models across five benchmark families and eight baselines, including
defense-aware adaptive attackers, with code, logs and generated tables released.
\end{enumerate}

% ============================================================================================
\section{Background and threat model}\label{sec:threat}
\subsection{Pipelines}
A generative application receives a user prompt $q$ and produces an output with a language model $f$. A
retrieval-augmented (RAG) application first retrieves the top-$k$ passages $c_1,\dots,c_k$ from a corpus
$\mathcal{C}$ and conditions $f$ on them~\cite{gao2023rag}. An agent interleaves plan steps $a_t$ (text and
proposed tool calls) with tool outputs $o_t$ for $t=1,\dots,T$. Every $q$, $c_i$ and $o_t$ is untrusted text
that reaches the model's context; every $a_t$ is the model's reaction to it.

\subsection{Taxonomy}
\begin{table}[t]\centering\small
\caption{Threat taxonomy. Evaluated classes are covered by Sections~\ref{sec:setup}--\ref{sec:results}.}
\label{tab:taxonomy}
\begin{tabular}{llllll}\toprule
\# & Threat class & Enters through & Property & Evaluated & Refs \\\midrule
T1 & Optimization-based jailbreak & GenAI prompt & Safety/integrity & Yes & \cite{zou2023gcg,mazeika2024harmbench}\\
T2 & Hand-crafted / in-the-wild jailbreaks & GenAI prompt & Safety & Yes & \cite{wei2023jailbroken,shen2024dan}\\
T3 & Direct prompt injection & GenAI prompt & Integrity/confid. & Yes & \cite{liu2023formalizing}\\
T4 & Indirect prompt injection & RAG context, observation & Integrity & Yes & \cite{greshake2023indirect,hines2024spotlighting}\\
T5 & Knowledge-base poisoning & RAG corpus & Integrity & Yes & \cite{zhong2023poisoning,zou2024poisonedrag}\\
T6 & Trigger-activated backdoor passages & RAG corpus & Integrity & Discussed & \cite{chaudhari2024phantom,xue2024badrag}\\
T7 & Jamming with blocker documents & RAG corpus & Availability & Yes & \cite{shafran2024jamming}\\
T8 & Datastore privacy leakage & RAG datastore & Confidentiality & Discussed & \cite{zeng2024goodbad}\\
T9 & Tool-call hijacking & Agent observation & Integrity & Yes & \cite{debenedetti2024agentdojo,zhan2024injecagent}\\
T10 & Exfiltration via improper tool use & Agent action & Confidentiality & Yes & \cite{zhan2024injecagent,fu2024imprompter}\\
T11 & Self-propagating injection & RAG + agent & Avail./integrity & Discussed & \cite{cohen2024worm}\\
T12 & Training-data poisoning; harmful agent compliance & Training, agent task & Integrity/safety & Discussed & \cite{carlini2024poisoning,andriushchenko2024agentharm,ruan2023toolemu,zhang2024asb}\\
\bottomrule\end{tabular}\end{table}

\subsection{Attacker and defender}
We consider three attacker profiles. A1 (black-box) can write to the prompt, corpus, web pages or tool
outputs but knows nothing about the defense. A2 (grey-box) additionally knows the public encoder, the
target false-alarm rate $\alpha$ and that WTS is deployed. A3 (white-box) also has the weights of the
protected open model and the layer used by the defense. The defender holds 50--1{,}000 benign samples per
channel, either the protected model's hidden states or an external encoder, and can drop, quarantine or
ask for confirmation, but cannot retrain the model; training-time defenses such as StruQ~\cite{chen2024struq}
and the instruction hierarchy~\cite{wallace2024hierarchy} are complementary.

% ============================================================================================
\section{Related work}\label{sec:related}
\TODO{1,000 words, comparative: attacks (T1--T12 above); defenses grouped as training-time
\cite{chen2024struq,wallace2024hierarchy}, prompt-level \cite{hines2024spotlighting}, input perturbation
\cite{robey2023smoothllm}, detection \cite{alon2023perplexity,inan2023llamaguard,liu2023formalizing,zou2025pishield},
RAG-specific \cite{xiang2024robustrag,choudhary2025stealth}; surveys \cite{yao2024survey,das2025security,deng2025agents};
optimal transport in ML \cite{montesuma2025ot}; conformal prediction \cite{angelopoulos2023conformal}.
State what each assumes that WTS does not, and the reverse.}

% ============================================================================================
\section{The Wasserstein Trust Score}\label{sec:method}
\subsection{Text units as distributions}
A frozen encoder $E$ maps a text unit $x$ to token vectors $h_1,\dots,h_n\in\mathbb{R}^d$, giving the
empirical measure $\mu_x=\frac1n\sum_{i=1}^n\delta_{h_i}$. A benign reference $\mu_{\mathrm{ref}}$ is built
in the same way from the pooled tokens of benign units of the same channel. Treating a unit as a cloud
rather than one pooled vector lets the score see a change in spread: an instruction appended to a benign
document can move the covariance of the cloud while barely moving its mean.

\subsection{Distance}
The Wasserstein-2 distance between measures $\mu$ and $\nu$ is
\begin{equation}
W_2^2(\mu,\nu)=\inf_{\pi\in\Pi(\mu,\nu)}\int\|a-b\|_2^2\,d\pi(a,b).
\end{equation}
Fitting $\mathcal{N}(m_x,\Sigma_x)$ to $\mu_x$ and $\mathcal{N}(m_r,\Sigma_r)$ to $\mu_{\mathrm{ref}}$ (Ledoit--Wolf
shrinkage) gives the closed form
\begin{equation}\label{eq:bw}
W_2^2=\underbrace{\|m_x-m_r\|_2^2}_{\text{location}}
+\underbrace{\tr\!\Big(\Sigma_x+\Sigma_r-2\big(\Sigma_r^{1/2}\Sigma_x\Sigma_r^{1/2}\big)^{1/2}\Big)}_{\text{shape}},
\end{equation}
which needs one symmetric eigendecomposition per unit because $\Sigma_r^{1/2}$ is computed once. The location
term is what centroid and Mahalanobis detectors see; the shape term is what they miss. As a check that
does not assume Gaussianity we also use the sliced distance
$SW_2^2(\mu,\nu)=\int_{S^{d-1}}W_2^2(\theta_\#\mu,\theta_\#\nu)\,d\sigma(\theta)$, estimated with $L=128$
directions. The Wasserstein Trust Score is $s(x)=W_2(\mu_x,\mu_{\mathrm{ref}})$.

\subsection{Three checkpoints}
\emph{Prompt (WTS-1).} $x=q$; the reference is benign prompts of the application.
\emph{Retrieved chunk (WTS-2).} For each $c_i$ in the top-$k$, the reference is the moment-matched mixture
$\lambda\,\mathcal{N}(m_{\mathcal{C}},\Sigma_{\mathcal{C}})+(1-\lambda)\,\mathcal{N}(m_{-i},\Sigma_{-i})$ of a
clean-corpus reference and the Gaussian of the other $k'-1$ retrieved chunks ($k'=20$, $\lambda=0.5$). A
flagged chunk is dropped and the context is refilled from lower ranks.
\emph{Agent (WTS-3).} Each tool output $o_t$ is scored against clean outputs of the same tool (a pooled model
is used when a tool has too few calibration outputs to certify $\alpha$). The intent drift
$D_t=W_2(\mu_{a_t},\mu_{a_{t-1}})$ compares the plan step produced after reading $o_t$ with the previous one.

\subsection{Calibration}
\begin{proposition}\label{prop:conformal}
Let $s_1,\dots,s_n$ be scores of benign calibration units disjoint from the reference data, and
$\tau_\alpha=s_{(\lceil(n+1)(1-\alpha)\rceil)}$. If a new benign unit is exchangeable with the calibration
units, then $\Pr\big(s(X_{n+1})>\tau_\alpha\big)\le\alpha$.
\end{proposition}
\begin{proof}
By exchangeability the rank of $s(X_{n+1})$ among $s_1,\dots,s_{n+1}$ is uniform on $\{1,\dots,n+1\}$ (ties
broken at random), so $\Pr(s(X_{n+1})>s_{(k)})\le (n+1-k)/(n+1)$ with $k=\lceil(n+1)(1-\alpha)\rceil$, which is at
most $\alpha$~\cite{angelopoulos2023conformal}.
\end{proof}
The guarantee holds for any score and any $n$ but needs $n\ge\lceil 1/\alpha\rceil-1$ for a finite threshold.
The agent guard combines the output and drift tests with conformal $p$-values and flags when
$\min(p_{\mathrm{out}},p_{\mathrm{drift}})\le\alpha/2$, so its false-alarm rate per step is at most $\alpha$ by
the union bound. Graded responses: \emph{block} (drop or redact), \emph{quarantine} (wrap as data) and
\emph{confirm} (withhold the proposed tool calls).

\subsection{Cost of evasion}
\begin{proposition}\label{prop:evasion}
Let the encoder be $L$-Lipschitz from the input metric to representation space under $W_2$. Any edit
$x\to x'$ with $s(x)>\tau_\alpha\ge s(x')$ satisfies $W_2(\mu_x,\mu_{x'})\ge s(x)-\tau_\alpha$, and hence moves
at least $(s(x)-\tau_\alpha)/L$ in the input metric.
\end{proposition}
\begin{proof}
By the triangle inequality, $s(x)=W_2(\mu_x,\mu_{\mathrm{ref}})\le W_2(\mu_x,\mu_{x'})+s(x')\le
W_2(\mu_x,\mu_{x'})+\tau_\alpha$; the Lipschitz bound converts the representation distance to the input metric.
\end{proof}
The bound is modest; Section~\ref{sec:adaptive} measures the actual cost of dilution attacks against it.

\subsection{Complexity}
Encoding dominates. Scoring costs one $d\times d$ eigendecomposition for the Gaussian score and
$O(Ln\log n)$ for the sliced score; Table~\ref{tab:e5} reports measured latency.

% ============================================================================================
\section{Experimental setup}\label{sec:setup}
\TODO{1,300 words from configs/e1--e4.yaml and environment.lock.md: models, encoders, benchmarks (pinned
commits), baselines, metrics, seeds, hardware, reproduction gates (reproduction.md).}

% ============================================================================================
\section{Results}\label{sec:results}
\subsection{Prompt channel (RQ1)}\IfFileExists{tables/table2_e1_detection.tex}{\input{tables/table2_e1_detection.tex}}{\TODO{run E1}}
\subsection{RAG (RQ2)}\IfFileExists{tables/table3_e2_rag.tex}{\input{tables/table3_e2_rag.tex}}{\TODO{run E2}}
\subsection{Agents (RQ3)}\IfFileExists{tables/table4_e3_agents.tex}{\input{tables/table4_e3_agents.tex}}{\TODO{run E3}}
\subsection{Calibration (RQ4)}\IfFileExists{tables/table_rq4_calibration.tex}{\input{tables/table_rq4_calibration.tex}}{\TODO{run E1}}
\subsection{Adaptive attacks (RQ5)}\label{sec:adaptive}\IfFileExists{tables/table5_e4_dilution.tex}{\input{tables/table5_e4_dilution.tex}}{\TODO{run E4}}
\subsection{Overhead (RQ6)}\IfFileExists{tables/table6_overhead.tex}{\input{tables/table6_overhead.tex}}{\TODO{run E5}\label{tab:e5}}
\subsection{Ablation of the score}\IfFileExists{tables/table7_ablation.tex}{\input{tables/table7_ablation.tex}}{\TODO{run E1}}
\TODO{One sentence per hypothesis H1--H6: supported / not supported, copied in substance from results/verdicts.md.}

\section{Discussion and limitations}\label{sec:discussion}
\TODO{Write the limitations from the failed or weak hypotheses first. Minimum: text only; English only;
stationarity of the benign distribution per channel; white-box results need hidden states; Gaussian fit on
very short units; no claim about attacks on the encoder itself.}

\section{Conclusion}\TODO{Two paragraphs; no new claims.}

\section*{Ethics statement}
All attacks were executed against public benchmarks and locally hosted or API-accessed models under their
terms of service; no deployed application or third party was targeted. Gated datasets are referenced by row
identifier and not redistributed.

\section*{Data and code availability}
\TODO{GitHub release tag and Zenodo DOI.}

\bibliographystyle{elsarticle-num}
\bibliography{refs}
\end{document}
